%! Planning a Study — Unit 3, Lesson 3.1 — Student Packet Cover Sheet
\documentclass[12pt]{article}
\usepackage{apstats-article}
\usepackage{apstats-boxes}
\usepackage{ltablex}
\keepXColumns

\begin{document}

% Full-bleed navy banner. \coverbanner MEASURES the title block and sizes the
% band to it, so a lesson title long enough to wrap grows the band rather than
% running out the bottom of it.
\coverbanner{Unit 3}{Lesson 3.1 \quad Planning a Study}

% NAMESTRIP: this is the ONLY component that carries the name row.
\namedateperiod
\vspace{0.18in}

\boxguard[14]
\begin{learningtargetbox}
\small
% GRADUAL RELEASE: the vocabulary is taught in the guided notes, so the targets name it
% plainly. The AP Learning Objectives (VAR-1.E, DAT-2.A, DAT-2.B) stay in the plan.
\begin{enumerate}[leftmargin=1.5em, itemsep=5pt]
  \item \ldots{} tell a \textbf{chance-based method} of collecting data from a
        \textbf{non-chance method}, and explain why only chance earns our trust.
  \item \ldots{} name the \textbf{population} and the \textbf{sample} in a study, say what a
        \textbf{census} is, and generalize only to the population that was sampled.
  \item \ldots{} classify a study as an \textbf{observational study} or an
        \textbf{experiment} by asking whether a treatment was imposed.
  \item \ldots{} explain why only a well-designed experiment can show \textbf{cause and
        effect}.
\end{enumerate}
\end{learningtargetbox}

\vspace{0.14in}

\boxguard[14]
\begin{tocbox}
\small
% A CONTENTS LIST ONLY: no Score column and no Total row. Whether a component counts is said in
% its description: AP Practice is extra credit, scored 1 to 5; the Homework is scored.
\renewcommand{\arraystretch}{1.5}
\begin{tabularx}{\linewidth}{@{} c l X @{}}
\rowcolor{sky}
\textbf{\#} & \textbf{Component} & \textbf{Description} \\
\midrule
1 & Warm-Up & An online poll about a later school start, a lunch survey, and sleep and test
                 scores \\
2 & Guided Notes \& Practice & Chance or people? Population, sample, and census; watching
                 versus experimenting --- then a university study app we work together \\
3 & AP Practice & A reading program, a state-park poll, and energy drinks and sleep: AP-style
                 multiple choice and one free-response set --- extra credit, scored 1 to 5 \\
4 & Homework & A city's dog-park survey and park lights: chance or not, the population and the
                 sample, and study or experiment --- \textbf{scored 1 to 5, due the first class
                 after two study halls} \\
\end{tabularx}
\end{tocbox}

\vspace{0.14in}

\boxguard[8]
\begin{remindbox}
\small
Before you believe a number, ask \textbf{how the data were collected}. If chance chose who is
in the data, no group gets a head start; if people chose, the result leans. A sample speaks only
for the \textbf{population} it was drawn from. And only an \textbf{experiment} --- where the
researchers \emph{impose} a treatment --- can show cause and effect. Measuring something very
carefully is still just watching.
\end{remindbox}

\end{document}
